The anchored interval of Equation~\eqref{eq:interval} has four parts: the family PD$(k)$ and its identified band, the second source that anchors $k$, the band of the anchored $k$ clipped to the identified set, and the outer Jeffreys limits. Table~\ref{tab:ablation} removes them one at a time at informativeness 1, where withdrawal is nearly uninformative, and at 9, where it is strongly informative, in the three designs of the coverage grid. The first row of each block is the baseline without any of the parts: the removed estimate with a Jeffreys interval.

\begin{table}[htbp]
\centering
\small
\caption{Ablation on grade B at informativeness 1 and 9: coverage (\%) in the first-year, pooled and ten-fold designs, and the width (percentage points) and relative bias of the first-year design, as the second source, the family PD$(k)$, the clipping to the identified set and the outer Jeffreys interval are retained or removed.}
\label{tab:ablation}
\begin{tabular}{p{0.9cm}p{6.1cm}rrrrr}
\toprule
 & Components retained & First & Pooled & $\times10$ & $W$ (pp) & Bias (\%) \\
\midrule
$\iota=1$ & Removed estimate with a Jeffreys interval (no anchor, no family) & 94.6 & 94.8 & 95.6 & 1.17 & $-0.2$ \\
 & Manski band (family and its bounds, no anchor) & 98.5 & 99.7 & 99.7 & 7.75 & -- \\
 & \relax Anchored $k$ band alone, mapped through PD$(k)$ and clipped to the identified set & 59.5 & 72.6 & 75.6 & 0.74 & $-0.0$ \\
 & Anchored interval (band, clipping and outer Jeffreys interval) & 98.5 & 99.7 & 99.6 & 1.99 & $-0.0$ \\
$\iota=9$ & Removed estimate with a Jeffreys interval (no anchor, no family) & 86.8 & 13.9 & 18.9 & 1.04 & $-20.8$ \\
 & Manski band (family and its bounds, no anchor) & 99.9 & 100.0 & 100.0 & 8.26 & -- \\
 & \relax Anchored $k$ band alone, mapped through PD$(k)$ and clipped to the identified set & 88.3 & 75.3 & 95.8 & 0.96 & $-0.1$ \\
 & Anchored interval (band, clipping and outer Jeffreys interval) & 99.9 & 99.2 & 100.0 & 2.16 & $-0.1$ \\
\bottomrule
\end{tabular}
\end{table}


The second source is what centres the estimate and narrows the interval. Without it the method keeps the family and the identified band, which covers \nv{cov.fy.i9.identified}\% at informativeness 9 in the first annual cohort and is \nv{w.fy.i9.identified} percentage points wide; the point that goes with it is the removed estimate, with its bias of $\nv{bias.fy.i9.industry}\%$. With the second source the point estimate has a bias of $\nv{bias.fy.i9.anchored}\%$ and the interval is \nv{w.fy.i9.anchored} points wide, \nv{w.x10.i9.anchored} in the cohort ten times larger. The second source therefore buys a reduction in width of about three quarters and the removal of the bias, on the condition of Proposition~\ref{prop:capture}: that it counts what the transition page would count.

The outer Jeffreys limits are what make the interval cover. The band of the anchored $k$ alone, which is the interval $[\pdk(k_{\ell}),\pdk(k_{u})]$ of the sensitivity analysis taken at face value, covers \nv{cov.fy.i9.band}\% at informativeness 9 in the first annual cohort and \nv{cov.fy.i1.band}\% at informativeness 1. The shortfall has a mechanism. The band counts the sampling noise of the second source and passes it through the family, but the family is clipped to the identified set: when the lower limit of the second source's rate is below the observed defaults $D/N$, the band's lower end is set to $D/N$, a value that the true PD is often below, because $D/N$ is itself a noisy count. Adding the Jeffreys limits of the transition page's own counts lowers that end to the Jeffreys limit of $D$ and raises the upper end in the same way, and the coverage rises to \nv{cov.fy.i9.anchored}\% and \nv{cov.fy.i1.anchored}\%. The price is the width: from \nv{w.fy.i9.band} to \nv{w.fy.i9.anchored} percentage points at informativeness 9. A narrower interval that covers at its nominal level would need a construction that treats the two sources' noise jointly rather than in sequence; this is not attempted here.

The clipping to the identified set and the family itself do not appear as separate rows because the engine applies them in every variant: the clipping keeps the point and the limits of the second source inside the interval that the counts allow, and the family is what turns a count from the second source into an odds ratio and a tipping value. The tipping $k$ adds nothing to the coverage; it adds the decision rule of Proposition~\ref{prop:tipping}, whose use the decision experiment of Table~\ref{tab:decisions} measures. The quality of the second source is the component with the largest effect on the centre of the result, and the mismatch ladder of Table~\ref{tab:capture} is its ablation: an exact source gives a bias of $\nv{bias.fy.i9.anchored}\%$, one that counts 0.8 of the defaults $\nv{cap.i9.c0.8.bias}\%$, and one that counts 1.3 times $+\nv{cap.i9.c1.3.abias}\%$, while the coverage stays above \nv{cap.i9.c1.3.cov}\% in all three.
